% Fixed texts, English (MSX2 version).
\usepackage[english]{babel}
\newcommand{\tManual}{User manual}
\newcommand{\tUserManual}{USER MANUAL}
\newcommand{\tForV}{for MSX2}
\newcommand{\tForTheV}{for MSX2}
\newcommand{\tCoverLine}{The Spectrum game, in colour on your MSX2}
\newcommand{\tOriginalBy}{An original game by}
\newcommand{\tPortBy}{MSX2 version by}
\newcommand{\tHomage}{This is an unofficial version made for study, preservation and as a homage.}
\newcommand{\tWhatYouNeed}{What you need}
\newcommand{\tItem}{Item}
\newcommand{\tRequirement}{Requirement}
\newcommand{\tComputer}{Computer}
\newcommand{\tComputerText}{\textbf{Any MSX2} (or MSX2+ or turbo R) with a free cartridge slot. It works at 50 and at 60~Hz, and the game runs at the same speed on both.}
\newcommand{\tMemory}{Memory}
\newcommand{\tMemoryText}{\textbf{\ramNeeded{} of RAM} and \textbf{128~KB of VRAM} (the usual amount on an MSX2).}
\newcommand{\tVideo}{Video}
\newcommand{\tVideoText}{The MSX2's own video chip (V9938 or V9958), in SCREEN~5. No video cartridge is needed.}
\newcommand{\tSoundRow}{Sound}
\newcommand{\tSoundText}{The MSX PSG: the music (title and game) and the sound effects.}
\newcommand{\tGameRow}{Game}
\newcommand{\tGameText}{The \textbf{\gameTitle{} cartridge for MSX2}: a \romKind{} (the file \texttt{\romFile}, written to a flash cartridge or similar).}
\newcommand{\tMonitor}{Monitor}
\newcommand{\tMonitorText}{The MSX's monitor or television.}
\newcommand{\tControlsRow}{Controls}
\newcommand{\tControlsText}{Cursor keys and the space bar, or a joystick in port 1 or port 2.}
\newcommand{\tNoTape}{\textbf{No tape, no disk, no password.} The whole game is in the cartridge.}
\newcommand{\tSettingUp}{Setting up}
\newcommand{\tSetOff}{\textbf{Switch the MSX off.} Never insert or remove a cartridge while the computer is on.}
\newcommand{\tSetGame}{Insert the \textbf{\gameTitle{} cartridge} in a cartridge slot.}
\newcommand{\tSetMonitor}{Connect the monitor or television as usual.}
\newcommand{\tSetJoy}{Plug a joystick into port 1 or port 2 if you want to use one.}
\newcommand{\tSetOn}{Switch the MSX on. After the computer's own start-up, the DiHalt Studio logo appears.}
\newcommand{\tBlackScreen}{\textbf{If the screen stays black}, check that the cartridge is properly inserted and that the computer is an MSX2 with \ramNeeded{} of RAM. In an emulator such as openMSX, load the ROM with the \textbf{ASCII16} mapper (option \texttt{-romtype ASCII16}).}
\newcommand{\tFlashCart}{\textbf{With a MegaFlashROM SCC+ SD:} copy the contents of the \texttt{dist} folder to the root of the SD card (the game as \texttt{ANDYCAPP.ROM}, the \texttt{OPFXSD} program, the files to start from the card and an \texttt{AUTOEXEC.BAT}) and switch the MSX on: the game is written to the cartridge in a few seconds and starts. From then on the cartridge starts \gameTitle{} every time you switch on. To start from the card, hold \key{ESC} while switching on (also the first time, if the cartridge already had another game written in it); \texttt{OPFXSD /C} erases the game from the cartridge.}
\newcommand{\tPowerOnTitle}{From power-on to the game}
\newcommand{\tPowerOnCaption}{Left: the cover screen. Right: the title screen.}
\newcommand{\tControlsTitle}{Controls}
\newcommand{\tCreditsTitle}{9. Credits and rights}
\newcommand{\tIsAGameBy}{is a game by}
\newcommand{\tRightsText}{Programmed by Blitter Animations; music by Jason C.~Brooke. The Andy Capp characters are creations of Reg Smythe and the property of Mirror Group Newspapers. The game, its graphics, music and name belong to their owners. This is an unofficial version made for study, preservation and as a homage; it is not affiliated with them.}
\newcommand{\tVersionWord}{MSX2 version}
\newcommand{\tMsxTrademark}{MSX is a trademark of MSX Licensing Corporation. The MSX logo is shown only to say which computers the game runs on}
\newcommand{\tSourceNote}{Source code available on GitHub. \textbf{Please improve the game if you feel so!}}
\newcommand{\tAppendix}{Appendix. Cartridge label}
\newcommand{\tStickerText}{A label for the cartridge, ready to print. At 100\,\% it is about 16\,cm wide; scale it in your printing program to the label area of your cartridge (the picture is 3:2).}
\newcommand{\tAction}{Action}
\newcommand{\tKeyboard}{Keyboard}
\newcommand{\tJoystick}{Joystick}
\newcommand{\tNo}{No.}
\newcommand{\tElement}{Element}
\newcommand{\tTellsYou}{What it tells you}
\newcommand{\tTips}{Tips.}
